% --- Archon physics formalization source begin ---
% archon:physics
% archon:covers PhyXMiniProblems/problem_phyx_mini_0882.lean
% archon:source-report reports/phyx_mini/problem_phyx_mini_0882.source.json

\chapter{Physics problem phyx\_mini\_0882}
\label{ch:PhyXMiniProblems_problem_phyx_mini_0882}

\paragraph{Problem source.}
\#\# Physical scenario

The switch in figure has been open for a long time. It is closed at \textbackslash{}( t = 0 \textbackslash{}, \textbackslash{}text\{s\} \textbackslash{}).

\#\# Question

After the switch has been closed for a long time, what is the current in the circuit?

\#\# Answer choices

- A: A: \textbackslash{}( \textbackslash{}Delta V\_\{\textbackslash{}text\{bat\}\}/2R \textbackslash{})

- B: B: \textbackslash{}( \textbackslash{}2Delta V\_\{\textbackslash{}text\{bat\}\}/R \textbackslash{})

- C: C: \textbackslash{}( \textbackslash{}Delta V\_\{\textbackslash{}text\{bat\}\}/R \textbackslash{})

- D: D: \textbackslash{}( \textbackslash{}Delta V\_\{\textbackslash{}text\{bat\}\}/3R \textbackslash{})

\#\# Figure caption (auxiliary; use the image as primary evidence)

The image depicts an electrical circuit diagram. The components and their relationships are as follows:

1. **Battery (\textbackslash{}( \textbackslash{}Delta V\_\{\textbackslash{}text\{bat\}\} \textbackslash{}))**: Located on the left side, symbolized by two parallel lines of different lengths, denoting the positive and negative terminals.

2. **Switch**: Positioned directly above the battery, it's shown in an open position and is capable of connecting two points of the circuit when closed.

3. **Resistor (\textbackslash{}( R \textbackslash{}))**: Placed to the right of the switch, depicted by a zigzag line, representing a resistive component in the circuit.

4. **Inductor (\textbackslash{}( L \textbackslash{}))**: Located to the far right, symbolized by a series of loops or a coil, indicating inductance within the circuit.

The circuit is arranged in series with the components connected in a loop, allowing current flow from the battery, through the switch, resistor, and inductor if the switch is closed.

\#\# Dataset metadata

Electric Circuits; Physical Model Grounding Reasoning, Implicit Condition Reasoning

\paragraph{Recorded answer/context.}
C: C: \textbackslash{}( \textbackslash{}Delta V\_\{\textbackslash{}text\{bat\}\}/R \textbackslash{})

\paragraph{Figure/image path.}
/root/proposal\_for\_physic/science-mango/phyx\_mini\_run/phyx\_data/test\_image/882.png

\paragraph{Formalization target.}
create a compiling Lean file with sorry bodies at `PhyXMiniProblems/problem\_phyx\_mini\_0882.lean`. The Lean declarations must preserve the physical quantities, dimensions or dimensional roles, figure labels, governing-law hypotheses, and final relation expressed by this problem.
Use Mathlib/Physlib names found through LeanExplore where available. If a physics API is missing, introduce faithful local abstractions rather than scalar placeholder aliases.

\begin{theorem}[Physics formalization target]
\label{thm:physics:phyx_mini_0882:target}
\lean{PhyXMiniProblems.ProblemPhyXMini0882.longTimeCurrent_eq_batteryVoltage_div_resistance}
\uses{def:physics:phyx-mini-0882:phyxminiproblems-problemphyxmini0882-currentinamperes, def:physics:phyx-mini-0882:phyxminiproblems-problemphyxmini0882-potentialdifferenceinvolts, def:physics:phyx-mini-0882:phyxminiproblems-problemphyxmini0882-resistanceinohms, def:physics:phyx-mini-0882:phyxminiproblems-problemphyxmini0882-seriesrlcircuitsetup, def:physics:phyx-mini-0882:phyxminiproblems-problemphyxmini0882-matchessuppliedseriesrlfigure, def:physics:phyx-mini-0882:phyxminiproblems-problemphyxmini0882-matchesseriesrltopology, def:physics:phyx-mini-0882:phyxminiproblems-problemphyxmini0882-matchesclosingexperiment, def:physics:phyx-mini-0882:phyxminiproblems-problemphyxmini0882-satisfiesidealseriesrldynamics, def:physics:phyx-mini-0882:phyxminiproblems-problemphyxmini0882-haslongtimeseriesrlsteadystate, def:physics:phyx-mini-0882:phyxminiproblems-problemphyxmini0882-displayedcurrentinamperes}
In the long-time steady state of the ideal series \(RL\) circuit, the current
is the battery voltage divided by the resistance.  This equals displayed
choice C, and positivity of the voltage and resistance makes C the unique
choice with that value.
\end{theorem}
\begin{proof}
For all sufficiently large times the switch is closed, so Ohm's law, the
inductor law, and Kirchhoff's loop law apply.  Pass to the limit in these
three identities.  The current tends to the stored long-time current, while
its derivative tends to zero.  Hence the resistor drop tends to
\(R I_\infty\), the inductor drop tends to zero, and Kirchhoff's law yields
\(\Delta V_{\rm bat}=R I_\infty\).  Since \(R>0\), division gives
\(I_\infty=\Delta V_{\rm bat}/R\).  The definition of displayed choice C is
exactly this quotient.  Finally inspect the four displayed formulas: after
multiplying by the positive resistance, equality with A, B, or D would force
the positive battery voltage to equal respectively one half, twice, or one
third of itself.  Each is impossible, so any matching choice is C.
\end{proof}

% --- Archon declaration topology begin ---
\section{Declaration topology}
\begin{definition}[electric Current Dimension]
  \label{def:physics:phyx-mini-0882:phyxminiproblems-problemphyxmini0882-electriccurrentdimension}
  \lean{PhyXMiniProblems.ProblemPhyXMini0882.electricCurrentDimension}
  The dimension C T⁻¹ of electric current
\end{definition}
\begin{proof}
  This is the declared model definition of electric Current Dimension.
\end{proof}
\begin{definition}[potential Difference Dimension]
  \label{def:physics:phyx-mini-0882:phyxminiproblems-problemphyxmini0882-potentialdifferencedimension}
  \lean{PhyXMiniProblems.ProblemPhyXMini0882.potentialDifferenceDimension}
  The dimension M L² T⁻² C⁻¹ of electric potential difference
\end{definition}
\begin{proof}
  This is the declared model definition of potential Difference Dimension.
\end{proof}
\begin{definition}[resistance Dimension]
  \label{def:physics:phyx-mini-0882:phyxminiproblems-problemphyxmini0882-resistancedimension}
  \lean{PhyXMiniProblems.ProblemPhyXMini0882.resistanceDimension}
  \uses{def:physics:phyx-mini-0882:phyxminiproblems-problemphyxmini0882-electriccurrentdimension, def:physics:phyx-mini-0882:phyxminiproblems-problemphyxmini0882-potentialdifferencedimension}
  The dimension V I⁻¹ of electrical resistance
\end{definition}
\begin{proof}
  This is the declared model definition of resistance Dimension.
\end{proof}
\begin{definition}[inductance Dimension]
  \label{def:physics:phyx-mini-0882:phyxminiproblems-problemphyxmini0882-inductancedimension}
  \lean{PhyXMiniProblems.ProblemPhyXMini0882.inductanceDimension}
  \uses{def:physics:phyx-mini-0882:phyxminiproblems-problemphyxmini0882-electriccurrentdimension, def:physics:phyx-mini-0882:phyxminiproblems-problemphyxmini0882-potentialdifferencedimension}
  The dimension V T I⁻¹ of electrical inductance
\end{definition}
\begin{proof}
  This is the declared model definition of inductance Dimension.
\end{proof}
\begin{definition}[Electric Current Quantity]
  \label{def:physics:phyx-mini-0882:phyxminiproblems-problemphyxmini0882-electriccurrentquantity}
  \lean{PhyXMiniProblems.ProblemPhyXMini0882.ElectricCurrentQuantity}
  \uses{def:physics:phyx-mini-0882:phyxminiproblems-problemphyxmini0882-electriccurrentdimension}
  A signed, unit-independent electric current
\end{definition}
\begin{proof}
  This is the declared model definition of Electric Current Quantity.
\end{proof}
\begin{definition}[Potential Difference Quantity]
  \label{def:physics:phyx-mini-0882:phyxminiproblems-problemphyxmini0882-potentialdifferencequantity}
  \lean{PhyXMiniProblems.ProblemPhyXMini0882.PotentialDifferenceQuantity}
  \uses{def:physics:phyx-mini-0882:phyxminiproblems-problemphyxmini0882-potentialdifferencedimension}
  A signed, unit-independent electric potential difference
\end{definition}
\begin{proof}
  This is the declared model definition of Potential Difference Quantity.
\end{proof}
\begin{definition}[Resistance Quantity]
  \label{def:physics:phyx-mini-0882:phyxminiproblems-problemphyxmini0882-resistancequantity}
  \lean{PhyXMiniProblems.ProblemPhyXMini0882.ResistanceQuantity}
  \uses{def:physics:phyx-mini-0882:phyxminiproblems-problemphyxmini0882-resistancedimension}
  A nonnegative, unit-independent electrical resistance
\end{definition}
\begin{proof}
  This is the declared model definition of Resistance Quantity.
\end{proof}
\begin{definition}[Inductance Quantity]
  \label{def:physics:phyx-mini-0882:phyxminiproblems-problemphyxmini0882-inductancequantity}
  \lean{PhyXMiniProblems.ProblemPhyXMini0882.InductanceQuantity}
  \uses{def:physics:phyx-mini-0882:phyxminiproblems-problemphyxmini0882-inductancedimension}
  A nonnegative, unit-independent electrical inductance
\end{definition}
\begin{proof}
  This is the declared model definition of Inductance Quantity.
\end{proof}
\begin{definition}[current In Amperes]
  \label{def:physics:phyx-mini-0882:phyxminiproblems-problemphyxmini0882-currentinamperes}
  \lean{PhyXMiniProblems.ProblemPhyXMini0882.currentInAmperes}
  \uses{def:physics:phyx-mini-0882:phyxminiproblems-problemphyxmini0882-electriccurrentquantity}
  Read a signed current in coherent-SI amperes
\end{definition}
\begin{proof}
  This is the declared model definition of current In Amperes.
\end{proof}
\begin{definition}[potential Difference In Volts]
  \label{def:physics:phyx-mini-0882:phyxminiproblems-problemphyxmini0882-potentialdifferenceinvolts}
  \lean{PhyXMiniProblems.ProblemPhyXMini0882.potentialDifferenceInVolts}
  \uses{def:physics:phyx-mini-0882:phyxminiproblems-problemphyxmini0882-potentialdifferencequantity}
  Read a signed potential difference in coherent-SI volts
\end{definition}
\begin{proof}
  This is the declared model definition of potential Difference In Volts.
\end{proof}
\begin{definition}[resistance In Ohms]
  \label{def:physics:phyx-mini-0882:phyxminiproblems-problemphyxmini0882-resistanceinohms}
  \lean{PhyXMiniProblems.ProblemPhyXMini0882.resistanceInOhms}
  \uses{def:physics:phyx-mini-0882:phyxminiproblems-problemphyxmini0882-resistancequantity}
  Read a nonnegative resistance in coherent-SI ohms
\end{definition}
\begin{proof}
  This is the declared model definition of resistance In Ohms.
\end{proof}
\begin{definition}[inductance In Henries]
  \label{def:physics:phyx-mini-0882:phyxminiproblems-problemphyxmini0882-inductanceinhenries}
  \lean{PhyXMiniProblems.ProblemPhyXMini0882.inductanceInHenries}
  \uses{def:physics:phyx-mini-0882:phyxminiproblems-problemphyxmini0882-inductancequantity}
  Read a nonnegative inductance in coherent-SI henries
\end{definition}
\begin{proof}
  This is the declared model definition of inductance In Henries.
\end{proof}
\begin{definition}[Circuit Component]
  \label{def:physics:phyx-mini-0882:phyxminiproblems-problemphyxmini0882-circuitcomponent}
  \lean{PhyXMiniProblems.ProblemPhyXMini0882.CircuitComponent}
  The four physical components visible in the single series loop
\end{definition}
\begin{proof}
  This is the declared model definition of Circuit Component.
\end{proof}
\begin{definition}[Circuit Node]
  \label{def:physics:phyx-mini-0882:phyxminiproblems-problemphyxmini0882-circuitnode}
  \lean{PhyXMiniProblems.ProblemPhyXMini0882.CircuitNode}
  Electrically distinct nodes encountered while traversing the loop
\end{definition}
\begin{proof}
  This is the declared model definition of Circuit Node.
\end{proof}
\begin{definition}[Switch Position]
  \label{def:physics:phyx-mini-0882:phyxminiproblems-problemphyxmini0882-switchposition}
  \lean{PhyXMiniProblems.ProblemPhyXMini0882.SwitchPosition}
  The two ideal states of the switch
\end{definition}
\begin{proof}
  This is the declared model definition of Switch Position.
\end{proof}
\begin{definition}[Battery Terminal]
  \label{def:physics:phyx-mini-0882:phyxminiproblems-problemphyxmini0882-batteryterminal}
  \lean{PhyXMiniProblems.ProblemPhyXMini0882.BatteryTerminal}
  The marked terminals of the battery symbol
\end{definition}
\begin{proof}
  This is the declared model definition of Battery Terminal.
\end{proof}
\begin{definition}[Series RLCircuit Figure]
  \label{def:physics:phyx-mini-0882:phyxminiproblems-problemphyxmini0882-seriesrlcircuitfigure}
  \lean{PhyXMiniProblems.ProblemPhyXMini0882.SeriesRLCircuitFigure}
  \uses{def:physics:phyx-mini-0882:phyxminiproblems-problemphyxmini0882-circuitcomponent, def:physics:phyx-mini-0882:phyxminiproblems-problemphyxmini0882-switchposition}
  Literal qualitative and text data transcribed from image 882.png. Printed labels are kept separate from the dimensionful component values
\end{definition}
\begin{proof}
  This is the declared model definition of Series RLCircuit Figure.
\end{proof}
\begin{definition}[Series RLCircuit Setup]
  \label{def:physics:phyx-mini-0882:phyxminiproblems-problemphyxmini0882-seriesrlcircuitsetup}
  \lean{PhyXMiniProblems.ProblemPhyXMini0882.SeriesRLCircuitSetup}
  \uses{def:physics:phyx-mini-0882:phyxminiproblems-problemphyxmini0882-electriccurrentquantity, def:physics:phyx-mini-0882:phyxminiproblems-problemphyxmini0882-potentialdifferencequantity, def:physics:phyx-mini-0882:phyxminiproblems-problemphyxmini0882-resistancequantity, def:physics:phyx-mini-0882:phyxminiproblems-problemphyxmini0882-inductancequantity, def:physics:phyx-mini-0882:phyxminiproblems-problemphyxmini0882-circuitcomponent, def:physics:phyx-mini-0882:phyxminiproblems-problemphyxmini0882-circuitnode, def:physics:phyx-mini-0882:phyxminiproblems-problemphyxmini0882-switchposition, def:physics:phyx-mini-0882:phyxminiproblems-problemphyxmini0882-batteryterminal, def:physics:phyx-mini-0882:phyxminiproblems-problemphyxmini0882-seriesrlcircuitfigure}
  Independent physical data for the switched series RL circuit. The two voltage-drop waveforms allow Ohm's law, the inductor constitutive law, and Kirchhoff's loop law to be stated separately. longTimeCurrent is not defined from the battery or resistance
\end{definition}
\begin{proof}
  This is the declared model definition of Series RLCircuit Setup.
\end{proof}
\begin{definition}[current Waveform In Amperes]
  \label{def:physics:phyx-mini-0882:phyxminiproblems-problemphyxmini0882-currentwaveforminamperes}
  \lean{PhyXMiniProblems.ProblemPhyXMini0882.currentWaveformInAmperes}
  \uses{def:physics:phyx-mini-0882:phyxminiproblems-problemphyxmini0882-currentinamperes, def:physics:phyx-mini-0882:phyxminiproblems-problemphyxmini0882-seriesrlcircuitsetup}
  The SI current readout as a function of Physlib time
\end{definition}
\begin{proof}
  This is the declared model definition of current Waveform In Amperes.
\end{proof}
\begin{definition}[current Derivative In Amperes Per Second]
  \label{def:physics:phyx-mini-0882:phyxminiproblems-problemphyxmini0882-currentderivativeinamperespersecond}
  \lean{PhyXMiniProblems.ProblemPhyXMini0882.currentDerivativeInAmperesPerSecond}
  \uses{def:physics:phyx-mini-0882:phyxminiproblems-problemphyxmini0882-seriesrlcircuitsetup, def:physics:phyx-mini-0882:phyxminiproblems-problemphyxmini0882-currentwaveforminamperes}
  The time derivative of the current readout, in amperes per second
\end{definition}
\begin{proof}
  This is the declared model definition of current Derivative In Amperes Per Second.
\end{proof}
\begin{definition}[Matches Supplied Series RLFigure]
  \label{def:physics:phyx-mini-0882:phyxminiproblems-problemphyxmini0882-matchessuppliedseriesrlfigure}
  \lean{PhyXMiniProblems.ProblemPhyXMini0882.MatchesSuppliedSeriesRLFigure}
  \uses{def:physics:phyx-mini-0882:phyxminiproblems-problemphyxmini0882-seriesrlcircuitsetup}
  The symbols, polarity marks, labels, and ordering visible in the raster
\end{definition}
\begin{proof}
  This is the declared model definition of Matches Supplied Series RLFigure.
\end{proof}
\begin{definition}[Matches Series RLTopology]
  \label{def:physics:phyx-mini-0882:phyxminiproblems-problemphyxmini0882-matchesseriesrltopology}
  \lean{PhyXMiniProblems.ProblemPhyXMini0882.MatchesSeriesRLTopology}
  \uses{def:physics:phyx-mini-0882:phyxminiproblems-problemphyxmini0882-seriesrlcircuitsetup}
  Node incidences expressing that all four components form one series loop
\end{definition}
\begin{proof}
  This is the declared model definition of Matches Series RLTopology.
\end{proof}
\begin{definition}[Matches Closing Experiment]
  \label{def:physics:phyx-mini-0882:phyxminiproblems-problemphyxmini0882-matchesclosingexperiment}
  \lean{PhyXMiniProblems.ProblemPhyXMini0882.MatchesClosingExperiment}
  \uses{def:physics:phyx-mini-0882:phyxminiproblems-problemphyxmini0882-seriesrlcircuitsetup}
  The switch has been open throughout negative time and is closed at t = 0 s and thereafter, exactly as stipulated in the problem statement
\end{definition}
\begin{proof}
  This is the declared model definition of Matches Closing Experiment.
\end{proof}
\begin{definition}[Satisfies Ideal Series RLDynamics]
  \label{def:physics:phyx-mini-0882:phyxminiproblems-problemphyxmini0882-satisfiesidealseriesrldynamics}
  \lean{PhyXMiniProblems.ProblemPhyXMini0882.SatisfiesIdealSeriesRLDynamics}
  \uses{def:physics:phyx-mini-0882:phyxminiproblems-problemphyxmini0882-potentialdifferenceinvolts, def:physics:phyx-mini-0882:phyxminiproblems-problemphyxmini0882-resistanceinohms, def:physics:phyx-mini-0882:phyxminiproblems-problemphyxmini0882-inductanceinhenries, def:physics:phyx-mini-0882:phyxminiproblems-problemphyxmini0882-seriesrlcircuitsetup, def:physics:phyx-mini-0882:phyxminiproblems-problemphyxmini0882-currentwaveforminamperes, def:physics:phyx-mini-0882:phyxminiproblems-problemphyxmini0882-currentderivativeinamperespersecond}
  Ideal lumped-element dynamics. These are governing laws, not the requested long-time formula: the resistor and inductor drops remain independent fields, and the battery/resistance quotient appears nowhere in this structure
\end{definition}
\begin{proof}
  This is the declared model definition of Satisfies Ideal Series RLDynamics.
\end{proof}
\begin{definition}[Has Long Time Series RLSteady State]
  \label{def:physics:phyx-mini-0882:phyxminiproblems-problemphyxmini0882-haslongtimeseriesrlsteadystate}
  \lean{PhyXMiniProblems.ProblemPhyXMini0882.HasLongTimeSeriesRLSteadyState}
  \uses{def:physics:phyx-mini-0882:phyxminiproblems-problemphyxmini0882-currentinamperes, def:physics:phyx-mini-0882:phyxminiproblems-problemphyxmini0882-seriesrlcircuitsetup, def:physics:phyx-mini-0882:phyxminiproblems-problemphyxmini0882-currentwaveforminamperes, def:physics:phyx-mini-0882:phyxminiproblems-problemphyxmini0882-currentderivativeinamperespersecond}
  The mathematical meaning of "after the switch has been closed for a long time": current approaches the independently stored steady current and its time derivative approaches zero. No value for that steady current is assumed
\end{definition}
\begin{proof}
  This is the declared model definition of Has Long Time Series RLSteady State.
\end{proof}
\begin{definition}[Answer Choice]
  \label{def:physics:phyx-mini-0882:phyxminiproblems-problemphyxmini0882-answerchoice}
  \lean{PhyXMiniProblems.ProblemPhyXMini0882.AnswerChoice}
  Labels of the four answer choices in the source
\end{definition}
\begin{proof}
  This is the declared model definition of Answer Choice.
\end{proof}
\begin{definition}[displayed Current In Amperes]
  \label{def:physics:phyx-mini-0882:phyxminiproblems-problemphyxmini0882-displayedcurrentinamperes}
  \lean{PhyXMiniProblems.ProblemPhyXMini0882.displayedCurrentInAmperes}
  \uses{def:physics:phyx-mini-0882:phyxminiproblems-problemphyxmini0882-potentialdifferenceinvolts, def:physics:phyx-mini-0882:phyxminiproblems-problemphyxmini0882-resistanceinohms, def:physics:phyx-mini-0882:phyxminiproblems-problemphyxmini0882-seriesrlcircuitsetup, def:physics:phyx-mini-0882:phyxminiproblems-problemphyxmini0882-answerchoice}
  The four displayed current readouts, interpreted in amperes
\end{definition}
\begin{proof}
  This is the declared model definition of displayed Current In Amperes.
\end{proof}
% --- Archon declaration topology end ---
% --- Archon physics formalization source end ---
